% Conclusion and Future Work
% Numbers come from reports/results_table.md, reports/ablation.csv, reports/checks and the Results section.

\section{Conclusion and Future Work}
\label{sec:conclusion}

\subsection{Summary}
We predicted next-day wildfire spread on the NDWS dataset with per-pixel gradient-boosted trees and
engineered spatial features. The tuned XGBoost reaches a test PR AUC of 0.257 and a ROC AUC of 0.906,
against 0.145-0.157 for models trained on the raw channels and 0.106 for persistence. At a 5\% false
positive rate it finds 68\% of next-day fire. LightGBM scores 0.248, so the result depends on the
features more than on the boosting library. Our score lies between the neighbour-feature random forest
of Huot et al.\ (0.225) and their CNN (0.284), and below recent U-Nets on the same split (0.346 in a
preprint~\cite{esparza2026}). Since these models come from different pipelines, this places our result
rather than ranking it, and our 95\% bootstrap interval over fire events (0.188-0.325) contains both of
Huot et al.'s values.

Three findings explain this result. First, the spatial features carry the signal. With fixed
hyperparameters, removing the fire-neighbourhood group drops PR AUC from 0.248 to 0.048, and that group
alone already reaches 0.218. Twenty random configurations per model changed eval PR AUC by only about
0.01, and training on 20 times less data gives almost the same eval score, so the remaining error lies
in the input representation rather than in tuning or data size. Second, the model behaves
physically: risk falls with distance to fire, rises with dry air and strong wind, and is pushed
downwind. In the training data, next-day fire 5\,km downwind is up to 3.1 times more frequent than
5\,km upwind in strong wind (Results section). Third, correct evaluation matters at 1\% prevalence. PR AUC, thresholds chosen on eval,
and scores per fire event give a much more honest picture than accuracy or ROC AUC.

\subsection{Limitations}
A per-pixel model with fixed-radius summaries learns a halo around today's fire. Recall is 0.99 for
pixels already burning, 0.65 at 2\,km and zero beyond 5\,km, where 64\% of the missed fire lies. False
alarms gather around fires that die down the next day: 89\% lie within 2\,km of day-$t$ fire, and
nothing in the inputs says which fires will be contained. Scored by event, the model still beats every
baseline (mean PR AUC 0.187), but at the pooled F1 threshold it finds fire in only 48\% of events,
fewer than persistence (61\%). Its raw scores are shifted by subsampling and class weighting. Undoing the shift in closed form gives
probabilities close to the observed fire rate overall, but up to about 3 times too low for the lowest-risk
pixels.

Two limits concern the evaluation itself. Every model was trained once, with one seed, on one split.
Resampling fire events gives wide intervals (0.188-0.325 for XGBoost) because a few large fires hold most
of the fire pixels. Paired resamples still separate the models (Appendix~\ref{app:checks}), and a rebuild
on another machine reproduced every PR AUC within 0.002. The week-based split separates samples in time, but a long-lived fire can appear in
weeks assigned to different splits, and we have not checked how often this happens.

\subsection{Future work: Phase~2}
Phase~2 moves the spatial reasoning from hand-made features into the model. A convolutional
network such as U-Net~\cite{ronneberger2015} sees the whole patch, including the shape of the fire
front and the wind field, which is exactly what the per-pixel model lacks. Because such a network
computes its loss on every pixel, it must handle the full 90:1 imbalance, and focal
loss~\cite{lin2017} is a natural candidate. Further directions are multi-day inputs (the best model on
WildfireSpreadTS uses five input days~\cite{gerard2023}); decision thresholds chosen per event, so that
small fires are not ignored; and calibration checked per region and fire size, so that the
probabilities can be read as risk.
